\documentclass[11pt]{article}

\usepackage[utf8]{inputenc}
\usepackage[T1]{fontenc}
\usepackage{lmodern}
\usepackage{geometry}
\usepackage{amsmath,amssymb,amsfonts,amsthm,mathtools}
\usepackage{bm}
\usepackage{enumitem}
\usepackage{microtype}
\usepackage{hyperref}
\usepackage{titlesec}
\usepackage{booktabs}
\usepackage{array}
\usepackage{setspace}
\usepackage{mathrsfs}
\usepackage{physics}

\geometry{margin=1in}
\hypersetup{colorlinks=true, linkcolor=black, urlcolor=blue, citecolor=black}
\setlist[enumerate]{topsep=0.4em,itemsep=0.35em}
\setlist[itemize]{topsep=0.3em,itemsep=0.25em}
\titleformat{\section}{\large\bfseries}{\thesection.}{0.5em}{}
\titleformat{\subsection}{\normalsize\bfseries}{\thesubsection.}{0.5em}{}
\setstretch{1.06}

\newcommand{\Aether}{\AE{}ther}
\newcommand{\Aflow}{\AE{}ther-flow}
\newcommand{\TheoryName}{\Aflow{} Interpretation of Relativity}
\newcommand{\TheoryProgram}{\Aether{} / \Aflow{} framework}
\newcommand{\Sub}{\mathcal A}
\newcommand{\ObserverMap}{\Xi}
\newcommand{\BridgeMetric}{\mathfrak g}
\newcommand{\ObserverMetric}{g^{\mathrm{br}}}
\newcommand{\CandidateMetric}{g^{\mathrm{cand}}}
\newcommand{\CompletedMetric}{g^{\sharp}}
\newcommand{\BaseShift}{\Sigma^{\mathrm{IER}}}
\newcommand{\ResponseShift}{\Sigma^{\mathrm{resp}}}
\newcommand{\CompShift}{\Sigma^{\mathrm{cmp}}}
\newcommand{\BaseLift}{\widehat\Sigma^{\mathrm{IER}}}
\newcommand{\ResponseLift}{\widehat\Sigma^{\mathrm{resp}}}
\newcommand{\CompLift}{\widehat\Sigma^{\mathrm{cmp}}}
\newcommand{\ReLongFour}{\widetilde{\mathcal R}_{\mu\nu}^{(\chi,\ge 4)}}
\newcommand{\ResidualCore}{\mathcal V_{\mu\nu}^{\mathrm{res}}}
\newcommand{\ResidualDec}{\mathcal D_{\mu\nu}^{\mathrm{dec}}}
\newcommand{\ResidualLift}{\widehat{\mathcal V}^{\mathrm{res}}}
\newcommand{\SubEinLin}{\widehat{\mathcal E}}
\newcommand{\SubGaugeVec}{\widehat{\mathcal B}}
\newcommand{\SubGaugeBook}{\widehat{\mathcal G}}
\newcommand{\SubReducedEin}{\widehat{\mathcal L}}
\newcommand{\FreeProj}{\Pi_{\mathrm{free}}}
\newcommand{\GaugeAux}{Y}
\newcommand{\ModeAux}{Z}
\newcommand{\ReducedEin}{\mathcal L_{\ObserverMetric}}
\newcommand{\GaugeBook}{\mathcal G_{\ObserverMetric}}
\newcommand{\EinLin}{\mathcal E_{\ObserverMetric}}
\newcommand{\EinQuad}{\mathcal Q_{\ge 2}^{\mathrm{Ein}}}
\newcommand{\CompMix}{\mathcal Q_{\mu\nu}^{\mathrm{cmp,mix}}}
\newcommand{\CompQuad}{\mathcal Q_{\mu\nu}^{\mathrm{cmp}}}
\newcommand{\CompDec}{\mathcal D_{\mu\nu}^{\sharp}}
\newcommand{\DeltaTComp}{\Delta T_{\mu\nu}^{\mathrm{cmp}}}
\newcommand{\MatterAct}{S_{\mathrm{matter}}}
\newcommand{\CompClass}{\mathscr C_{\mathrm{cmp}}}
\newcommand{\CompFree}{\mathscr C_{\mathrm{cmp}}^{\mathrm{free}}}
\newcommand{\CompForced}{\mathscr C_{\mathrm{cmp}}^{\mathrm{for}}}
\newcommand{\CompKernel}{\mathscr K_{\mathrm{cmp}}}
\newcommand{\MicFunctional}{\mathscr F_{\mathrm{mic}}}
\newcommand{\FreeStiff}{\kappa_{\mathrm{free}}}
\newcommand{\epsIR}{\varepsilon}

\newtheorem{definition}{Definition}
\newtheorem{proposition}{Proposition}
\newtheorem{remark}{Remark}
\newtheorem{corollary}{Corollary}
\newtheorem{theorem}{Theorem}

\title{\textbf{The \TheoryName{}}\\[0.4em]
\large Substrate-Response / Self-Response Charge-Polarization Candidate\\
\large Quartic-Completion Selection-Principle Deeper-Origin Note}
\author{}
\date{}

\begin{document}

\maketitle

\begin{abstract}
The quartic-completion kernel-origin note already fixed, at disciplined current-scope selection-principle level, why the same completion kernel, reduced-gauge auxiliary block, and free-mode projector are selected on the one-metric charge-polarization branch. The remaining burden has therefore shifted one layer deeper again. It is no longer to write another same-level completion action. It is to explain why branch-preserving quadratic response, reduced-gauge neutrality, and equilibrium suppression of source-free homogeneous data arise in the first place from a deeper substrate response principle.

The present note gives that still-deeper origin at controlled microscopic-equilibrium scope. The key input is a neutral branch-equilibrium completion reservoir around the already-fixed charge-polarization branch. Three structural requirements are imposed. First, the unsourced branch is a stationary equilibrium of the same bridge geometry and observer map, so the first nontrivial completion response is the quadratic Hessian of the same branch-preserving bridge action in the completion direction. Second, microscopic weights are invariant under reduced-gauge relabeling of completion data and carry no net reduced-gauge charge, so the completion reservoir depends only on the neutral combination \(\GaugeAux-\SubGaugeVec(H)\), which forces the reduced-gauge auxiliary block and the same reduced operator \(\SubReducedEin=\SubEinLin-\SubGaugeBook\). Third, source-free homogeneous completion data carry positive equilibrium stiffness and the lifted quartic source \(\ResidualLift\) has no component in that free sector, so equilibrium forces the free-mode projector condition \(\FreeProj H=0\).

These deeper inputs yield the microscopic equilibrium functional
\[
\begin{aligned}
\MicFunctional[H,\GaugeAux]
=
\int_{\Sub}\sqrt{G}\,
\Bigl[
&\frac12 H^{AB}\SubEinLin(H)_{AB}
+\frac12\bigl(\GaugeAux_A-\SubGaugeVec(H)_A\bigr)
\bigl(\GaugeAux^A-\SubGaugeVec(H)^A\bigr)
\\
&+\frac{\FreeStiff}{2}(\FreeProj H)^{AB}(\FreeProj H)_{AB}
+(\ResidualLift)^{AB}H_{AB}
\Bigr],
\end{aligned}
\]
whose stationary point satisfies
\[
\SubReducedEin(\CompLift)_{AB}=-\ResidualLift_{AB},
\qquad
\FreeProj\CompLift=0.
\]
Hence the same observer solve
\[
\ReducedEin(\CompShift)_{\mu\nu}=-\ResidualCore{}_{\mu\nu}
\]
is recovered once more, together with the same one operative metric, the same Phase D / E infrared and universal-coupling package, the same Phase F controlled GR limit, and the same removal of the former genuine quartic residual tensor from the exact observer equation.

The result is positive but still disciplined. The note supplies a deeper microscopic origin of the current-scope selection principle. It does not claim a final ultraviolet completion of the branch, a uniqueness statement beyond the fixed local two-derivative scope, or a wider theorem than the present one-metric GR claim boundary allows.
\end{abstract}

\section{Introduction}

The active charge-polarization line has already crossed several exact milestones on the same one-metric GR route. The explicit candidate is fixed. The Phase D / E infrared-spectrum and universal-coupling package is fixed. The Phase F controlled GR-limit note is fixed. The residual-sector verdict is fixed. The quartic-completion redesign is fixed. The substrate-anchoring note is fixed. The substrate-action note is fixed. The kernel-origin note is fixed.

The honest next burden is therefore sharper again. It is not another tracking-document pass. It is not another same-layer completion-sector action note. It is not, by default, another deeper anchored positive-pair manuscript. The remaining primary burden on the GR line is now the still-deeper microscopic origin of the selection principle recorded in the kernel-origin note.

\paragraph{Framework and claim boundary.}
\TheoryProgram{} denotes the broader \Aether{} / \Aflow{} framework built on the claims that the \Aether{} is the underlying four-dimensional substrate of reality, the \Aflow{} is its intrinsic ordered motion, observed three-dimensional space is the local experiential slice of that deeper substrate, S-time is the experienced order of change arising from matter, light, and the \Aflow{}, and observed expansion is the three-dimensional appearance of deeper four-dimensional ordered motion. Within that framework, \TheoryName{} denotes the exact-closure theory in which the effective gravitational dynamics are adopted to be exactly Einsteinian with universal matter coupling. In the active sequence, \emph{adoption} means use of that established relativistic dynamics without claiming substrate derivation, while \emph{derivation} is reserved for a first-principles recovery from explicit substrate variables.

The present note stays strictly inside that boundary. It does not claim a final ultraviolet completion. It does not claim that every possible microscopic model is now classified. It does not widen the current finite-amplitude theorem boundary. Its narrower task is exactly the next one recorded in the claim-boundary and project-tracking Markdown: derive or justify why the already-selected completion kernel, reduced-gauge neutrality condition, and free-mode projector arise from deeper branch-equilibrium substrate structure while preserving the same one operative metric, the same Phase D / E package, the same controlled GR limit, and the same removal of the former genuine quartic residual sector.

\section{Fixed Inputs from the Recorded Completion Line}

\begin{definition}[Fixed charge-polarization branch and source]
\label{def:fixed}
The present note takes as fixed the same charge-polarization branch
\begin{equation}
\CandidateMetric
=
\ObserverMetric+\BaseShift+\ResponseShift,
\qquad
\CompletedMetric
=
\ObserverMetric+\BaseShift+\ResponseShift+\CompShift,
\label{eq:metrics}
\end{equation}
with lifted completion variables \(\BaseLift,\ResponseLift,\CompLift\) satisfying
\begin{equation}
\ObserverMap^*\BaseLift=\BaseShift,
\qquad
\ObserverMap^*\ResponseLift=\ResponseShift,
\qquad
\ObserverMap^*\CompLift=\CompShift.
\label{eq:lifts}
\end{equation}
The unresolved observer-side quartic tensor and its fixed substrate lift are
\begin{equation}
\ResidualCore
\equiv
\ReLongFour
+\EinQuad[\BaseShift]
+\EinLin(\ResponseShift),
\label{eq:vres}
\end{equation}
\begin{equation}
\ResidualLift_{AB}
\equiv
\widehat{\mathcal R}_{AB}^{(\chi,\ge 4)}
+\widehat{\mathcal Q}_{AB}^{\mathrm{Ein}}[\BaseLift]
+\widehat{\mathcal E}_{AB}(\ResponseLift),
\qquad
\ObserverMap^*\ResidualLift=\ResidualCore.
\label{eq:liftedsource}
\end{equation}
The fixed orders are
\begin{equation}
\BaseShift=\mathcal O_{\ge 2},
\qquad
\ResponseShift=\mathcal O_{\ge 4},
\qquad
\CompShift=\mathcal O_{\ge 4},
\qquad
\ResidualLift=\mathcal O_{\ge 4},
\label{eq:orders}
\end{equation}
and on every controlled infrared family,
\begin{equation}
\BaseShift=\mathcal O(\epsIR^2),
\qquad
\ResponseShift=\mathcal O(\epsIR^4),
\qquad
\CompShift=\mathcal O(\epsIR^4),
\qquad
\ResidualLift=\mathcal O(\epsIR^4).
\label{eq:irorders}
\end{equation}
\end{definition}

\begin{definition}[Fixed bridge-response operators]
\label{def:operators}
Let \(\widehat\nabla\) denote the Levi-Civita connection of the unshifted bridge metric \(\BridgeMetric\). For any observer-compatible symmetric substrate tensor \(H_{AB}\), define
\begin{equation}
\SubGaugeVec(H)_A
\equiv
\widehat\nabla^B\!\left(
H_{AB}
-\frac12 \BridgeMetric_{AB}\BridgeMetric^{CD}H_{CD}
\right),
\label{eq:subB}
\end{equation}
\begin{equation}
\SubGaugeBook(H)_{AB}
\equiv
\widehat\nabla_{(A}\SubGaugeVec(H)_{B)}
-\frac12 \BridgeMetric_{AB}\widehat\nabla^C\SubGaugeVec(H)_C,
\label{eq:subG}
\end{equation}
\begin{equation}
\SubReducedEin(H)_{AB}
\equiv
\SubEinLin(H)_{AB}-\SubGaugeBook(H)_{AB},
\label{eq:subL}
\end{equation}
chosen so that
\begin{equation}
\ObserverMap^*\bigl(\SubEinLin(H)\bigr)=\EinLin(\ObserverMap^*H),
\qquad
\ObserverMap^*\bigl(\SubGaugeBook(H)\bigr)=\GaugeBook(\ObserverMap^*H),
\label{eq:intertwine1}
\end{equation}
and hence
\begin{equation}
\ObserverMap^*\bigl(\SubReducedEin(H)\bigr)
=
\ReducedEin(\ObserverMap^*H).
\label{eq:intertwine2}
\end{equation}
\end{definition}

\begin{definition}[Admissible completion class and free sector]
\label{def:class}
Let \(\CompClass\) denote the same observer-compatible reduced-harmonic Coulomb-plus-regular completion class used in the redesign, anchoring, action, and kernel-origin notes. The free homogeneous sector is
\begin{equation}
\CompFree
\equiv
\left\{
H\in\CompClass:\SubReducedEin(H)=0
\right\},
\label{eq:freeclass}
\end{equation}
and \(\CompForced\) denotes its orthogonal complement inside \(\CompClass\) with respect to the bridge-metric \(L^2\) pairing
\begin{equation}
\langle H_1,H_2\rangle_{\mathrm{cmp}}
\equiv
\int_{\Sub} d^4X\,\sqrt{G}\,
\BridgeMetric^{AC}\BridgeMetric^{BD}(H_1)_{AB}(H_2)_{CD}.
\label{eq:inner}
\end{equation}
Let \(\FreeProj:\CompClass\to\CompFree\) be the orthogonal projector. The fixed source satisfies
\begin{equation}
\ResidualLift\in\CompForced,
\qquad
\FreeProj\ResidualLift=0.
\label{eq:sourcefree}
\end{equation}
\end{definition}

\begin{remark}
The present note does not alter the source \(\ResidualLift\), the reduced operator \(\SubReducedEin\), the one-metric matter route, or the recorded admissible completion class. It explains why the selection principles used in the kernel-origin note arise from deeper microscopic-equilibrium structure below that note's scope.
\end{remark}

\section{Microscopic-Equilibrium Origin Data}

\begin{definition}[Deeper branch-equilibrium completion origin]
\label{def:micro}
An \emph{admissible deeper origin of the current-scope completion selection principle} is a coarse-grained microscopic completion reservoir satisfying all of the following conditions.
\begin{enumerate}
    \item \textbf{Branch stationarity.} With the quartic source turned off, the reference branch \(H=0\) is a stationary equilibrium of the same bridge metric \(\BridgeMetric\), observer map \(\ObserverMap\), and one-metric matter route already fixed in the active sequence.
    \item \textbf{Local branch covariance.} At the first unresolved order, the completion reservoir may depend only on the fixed bridge data, on observer-compatible completion tensors \(H\in\CompClass\), and on the already-fixed source \(\ResidualLift\). No new metric route, no new direct matter coupling, and no new higher-derivative principal operator may be inserted.
    \item \textbf{Reduced-gauge relabeling neutrality.} Microscopic weights are invariant under reduced-gauge relabelings of completion data and carry no net reduced-gauge charge. Hence gauge-exact completion directions are neutral bookkeeping directions rather than new physical observer modes.
    \item \textbf{Positive free-mode stiffness.} The source-free homogeneous completion sector \(\CompFree\) carries strictly positive equilibrium stiffness. Unsourced free homogeneous data therefore cost positive branch free energy rather than representing flat directions of the physical branch.
    \item \textbf{Minimal source loading.} The fixed lifted quartic tensor \(\ResidualLift\) is the only current-scope loading source at first unresolved order.
\end{enumerate}
\end{definition}

\begin{remark}
Definition \ref{def:micro} is deliberately weaker than a final ultraviolet model. The object being justified here is not every microscopic coefficient in every possible completion. It is the selection principle itself: why the current branch responds quadratically, why reduced-gauge neutrality is forced, and why free homogeneous completion data are suppressed at equilibrium.
\end{remark}

\begin{proposition}[Branch stationarity forces quadratic response]
\label{prop:quadratic}
Under Definition \ref{def:micro}(1), the first nontrivial unsourced completion response about the reference branch is quadratic in the completion variable \(H\).
\end{proposition}

\begin{proof}
Let \(\MicFunctional^{(0)}\) denote the unsourced microscopic branch free energy. By branch stationarity, \(H=0\) is a stationary equilibrium, so the first variation of \(\MicFunctional^{(0)}\) at \(H=0\) vanishes. Therefore the Taylor expansion around the reference branch begins with the Hessian:
\[
\MicFunctional^{(0)}[H]
=
\MicFunctional^{(0)}[0]
+\frac12 \delta^2\MicFunctional^{(0)}[0](H,H)
+\mathcal O(H^3).
\]
Hence the first nontrivial response is quadratic in \(H\). This is the deeper origin of the kernel-origin note's statement that the first unresolved completion response is quadratic: it is the ordinary consequence of stationary branch equilibrium at the unsourced reference point.
\end{proof}

\begin{proposition}[Local branch covariance identifies the quadratic Hessian]
\label{prop:hessian}
Under Definition \ref{def:micro}(2), the quadratic \(H\)-H Hessian selected at the first unresolved order is the same branch-preserving bridge-response form
\begin{equation}
\frac12
\int_{\Sub} d^4X\,\sqrt{G}\,
H^{AB}\SubEinLin(H)_{AB}.
\label{eq:HHhessian}
\end{equation}
\end{proposition}

\begin{proof}
By Proposition \ref{prop:quadratic}, only the Hessian of the branch free energy enters at first nontrivial order. Definition \ref{def:micro}(2) forbids any new metric route, any new direct matter coupling, and any new higher-derivative principal operator at that stage. Therefore the Hessian must be built from the same local two-derivative bridge geometry already present on the fixed branch. The only already-recorded symmetric observer-compatible two-tensor response in that fixed bridge geometry is the bridge Einstein--Hilbert second variation \(\SubEinLin\). Any different first-unresolved-order Hessian would introduce an extra principal block not present on the fixed branch or would spoil the established intertwining with the observer-side linear Einstein operator. Hence \eqref{eq:HHhessian} is selected at the present disciplined microscopic scope.
\end{proof}

\begin{remark}
Proposition \ref{prop:hessian} is not a claim that every conceivable ultraviolet completion in principle must reduce to \(\SubEinLin\). The claim is narrower and exactly tuned to the active branch: once one keeps the same bridge data, the same observer map, the same derivative order, and the same one-metric matter route, the unsourced completion Hessian is forced to be the already-recorded bridge response rather than an arbitrary new insertion.
\end{remark}

\section{Reduced-Gauge Neutrality from Microscopic Relabeling Symmetry}

\begin{definition}[Neutral relabeling realization]
\label{def:neutral}
At current scope, a \emph{neutral relabeling realization} of reduced-gauge symmetry is a local quadratic microscopic term built from \(H\) and an auxiliary covector \(\GaugeAux_A\) such that:
\begin{enumerate}
    \item it is invariant under simultaneous reduced-gauge relabeling of \(H\) and \(\GaugeAux_A\);
    \item it introduces no new observer-accessible pole content;
    \item it is first-order in derivatives of \(H\);
    \item after elimination of \(\GaugeAux_A\), it produces only reduced-gauge bookkeeping on the observer side.
\end{enumerate}
\end{definition}

\begin{proposition}[Neutrality forces the combination \(\GaugeAux-\SubGaugeVec(H)\)]
\label{prop:neutralcombo}
Under Definition \ref{def:micro}(3), the unique current-scope quadratic neutral relabeling realization is
\begin{equation}
\frac12
\int_{\Sub} d^4X\,\sqrt{G}\,
\bigl(\GaugeAux_A-\SubGaugeVec(H)_A\bigr)
\bigl(\GaugeAux^A-\SubGaugeVec(H)^A\bigr).
\label{eq:neutralterm}
\end{equation}
\end{proposition}

\begin{proof}
Reduced-gauge relabeling neutrality requires the microscopic free energy to depend only on a neutral quantity built from the auxiliary covector and a first-derivative covector of \(H\). Definition \ref{def:neutral}(3) restricts that covector to first derivative order. Among current-scope observer-compatible first-derivative covectors, the fixed reduced-gauge vector \(\SubGaugeVec(H)\) is the one whose symmetrized derivative gives exactly the observer gauge-bookkeeping tensor under pullback. Any different choice would either fail to be the reduced-gauge relabeling current already fixed in Definition \ref{def:operators}, introduce new derivative structure, or generate new observer-side content. The unique quadratic neutral invariant is therefore the square of the neutral combination \(\GaugeAux-\SubGaugeVec(H)\), namely \eqref{eq:neutralterm}.
\end{proof}

\begin{proposition}[Neutrality reproduces the reduced operator]
\label{prop:reduced}
Combining \eqref{eq:HHhessian} with \eqref{eq:neutralterm} yields the reduced quadratic response
\begin{equation}
\frac12
\int_{\Sub} d^4X\,\sqrt{G}\,
H^{AB}\SubReducedEin(H)_{AB}
-\int_{\Sub} d^4X\,\sqrt{G}\,
\GaugeAux^A\SubGaugeVec(H)_A
+\frac12
\int_{\Sub} d^4X\,\sqrt{G}\,
\GaugeAux^A\GaugeAux_A.
\label{eq:reducedfunctional}
\end{equation}
In particular, reduced-gauge neutrality forces the same reduced completion operator \(\SubReducedEin=\SubEinLin-\SubGaugeBook\) used in the kernel-origin note.
\end{proposition}

\begin{proof}
Expand \eqref{eq:neutralterm}:
\[
\frac12\int_{\Sub}\sqrt{G}\,\GaugeAux^A\GaugeAux_A
-\int_{\Sub}\sqrt{G}\,\GaugeAux^A\SubGaugeVec(H)_A
+\frac12\int_{\Sub}\sqrt{G}\,\SubGaugeVec(H)^A\SubGaugeVec(H)_A.
\]
The last term is exactly the reduced-gauge bookkeeping correction. By integration by parts in the admissible completion class,
\[
\frac12\int_{\Sub}\sqrt{G}\,\SubGaugeVec(H)^A\SubGaugeVec(H)_A
=
-\frac12\int_{\Sub}\sqrt{G}\,H^{AB}\SubGaugeBook(H)_{AB},
\]
up to the already-fixed vanishing boundary contribution in the Coulomb-plus-regular class. Adding this to the Hessian term \(\frac12\int H^{AB}\SubEinLin(H)_{AB}\) gives \(\frac12\int H^{AB}\SubReducedEin(H)_{AB}\). Thus \eqref{eq:reducedfunctional} follows.
\end{proof}

\begin{corollary}[Deeper origin of reduced-gauge neutrality]
\label{cor:neutral}
Reduced-gauge neutrality is not merely imposed because it is convenient for the current-scope solve. It is forced by the deeper microscopic statement that the completion reservoir is neutral under reduced-gauge relabeling and can depend only on the neutral combination \(\GaugeAux-\SubGaugeVec(H)\).
\end{corollary}

\section{Free-Mode Suppression from Positive Equilibrium Stiffness}

\begin{definition}[Positive free-mode stiffness]
\label{def:stiff}
The microscopic equilibrium cost of source-free homogeneous completion data is
\begin{equation}
\frac{\FreeStiff}{2}
\int_{\Sub} d^4X\,\sqrt{G}\,
(\FreeProj H)^{AB}(\FreeProj H)_{AB},
\qquad
\FreeStiff>0.
\label{eq:stiffness}
\end{equation}
\end{definition}

\begin{remark}
Equation \eqref{eq:stiffness} is the deeper equilibrium statement lying below the kernel-origin note's zero-free-mode rule. The free homogeneous sector is not treated as another independent physical branch direction. It is treated as an unsourced equilibrium excitation with positive cost.
\end{remark}

\begin{proposition}[Positive stiffness forces \(\FreeProj H=0\)]
\label{prop:stiffzero}
Let \(H\in\CompClass\) be a stationary point of the microscopic completion functional built from \eqref{eq:HHhessian}, \eqref{eq:neutralterm}, \eqref{eq:stiffness}, and the minimal source loading \(\int \sqrt{G}\,(\ResidualLift)^{AB}H_{AB}\). Then
\begin{equation}
\FreeProj H=0.
\label{eq:freezero}
\end{equation}
\end{proposition}

\begin{proof}
After using Proposition \ref{prop:reduced}, stationarity in the \(H\)-direction gives an equation of the form
\begin{equation}
\SubReducedEin(H)_{AB}
+\FreeStiff(\FreeProj H)_{AB}
=
-\ResidualLift_{AB}.
\label{eq:preproject}
\end{equation}
Apply \(\FreeProj\) to \eqref{eq:preproject}. By Definition \ref{def:class}, \(\CompFree=\ker\SubReducedEin\). Since \(\SubReducedEin\) is symmetric with respect to \(\langle\cdot,\cdot\rangle_{\mathrm{cmp}}\), its range lies in \(\CompForced\), so \(\FreeProj\SubReducedEin(H)=0\). Equation \eqref{eq:sourcefree} gives \(\FreeProj\ResidualLift=0\). Therefore \eqref{eq:preproject} reduces to
\[
\FreeStiff\,\FreeProj H=0.
\]
Because \(\FreeStiff>0\), it follows that \(\FreeProj H=0\).
\end{proof}

\begin{corollary}[Deeper origin of the projector condition]
\label{cor:projector}
The kernel-origin note's equilibrium rule \(\FreeProj H=0\) arises from the deeper microscopic fact that the free homogeneous completion sector carries positive equilibrium stiffness while the lifted quartic source carries no free homogeneous component.
\end{corollary}

\section{Microscopic Completion Functional and Recovery of the Same Solve}

\begin{definition}[Microscopic completion functional]
\label{def:micfunc}
The current-scope deeper-origin microscopic completion functional is
\begin{equation}
\begin{aligned}
\MicFunctional[H,\GaugeAux]
\equiv
\int_{\Sub} d^4X\,\sqrt{G}\,
\Bigl[
&\frac12 H^{AB}\SubEinLin(H)_{AB}
+\frac12\bigl(\GaugeAux_A-\SubGaugeVec(H)_A\bigr)
\bigl(\GaugeAux^A-\SubGaugeVec(H)^A\bigr)
\\
&+\frac{\FreeStiff}{2}(\FreeProj H)^{AB}(\FreeProj H)_{AB}
+(\ResidualLift)^{AB}H_{AB}
\Bigr].
\end{aligned}
\label{eq:micfunc}
\end{equation}
\end{definition}

\begin{remark}
Definition \ref{def:micfunc} is not introduced as another same-level completion-sector action note. It is a still-deeper microscopic-equilibrium justification of the already-recorded selection principle. Its on-shell consequence will be exactly the same \(\Sigma^{\mathrm{cmp}}\) solve and the same one-metric observer branch.
\end{remark}

\begin{theorem}[The kernel-origin selection principles follow from \(\MicFunctional\)]
\label{thm:principles}
The microscopic functional \eqref{eq:micfunc} yields the three current-scope selection principles of the kernel-origin note:
\begin{enumerate}
    \item branch-preserving quadratic response arises from stationary branch equilibrium and the fixed bridge Hessian;
    \item reduced-gauge neutrality is forced by microscopic relabeling neutrality and yields the reduced operator \(\SubReducedEin\);
    \item equilibrium suppression of source-free homogeneous data is forced by the positive free-mode stiffness \(\FreeStiff\).
\end{enumerate}
\end{theorem}

\begin{proof}
Item (1) is Proposition \ref{prop:quadratic} and Proposition \ref{prop:hessian}. Item (2) is Proposition \ref{prop:neutralcombo}, Proposition \ref{prop:reduced}, and Corollary \ref{cor:neutral}. Item (3) is Definition \ref{def:stiff}, Proposition \ref{prop:stiffzero}, and Corollary \ref{cor:projector}.
\end{proof}

\begin{theorem}[Stationary equations]
\label{thm:stationary}
Let \(\CompLift\in\CompClass\) and \(\GaugeAux\) be a stationary point of \(\MicFunctional\). Then
\begin{equation}
\GaugeAux_A=\SubGaugeVec(\CompLift)_A,
\label{eq:Yeq}
\end{equation}
and
\begin{equation}
\SubReducedEin(\CompLift)_{AB}
+\FreeStiff(\FreeProj\CompLift)_{AB}
=
-\ResidualLift_{AB}.
\label{eq:Heq}
\end{equation}
\end{theorem}

\begin{proof}
Variation of \eqref{eq:micfunc} with respect to \(\GaugeAux_A\) gives \eqref{eq:Yeq}. Substituting that relation back into the quadratic neutral term reproduces the reduced functional of Proposition \ref{prop:reduced}. Variation with respect to \(H_{AB}\) then yields \eqref{eq:Heq}, since the derivative of the stiffness term \eqref{eq:stiffness} is \(\FreeStiff\,\FreeProj H\).
\end{proof}

\begin{corollary}[The same substrate completion solve]
\label{cor:subsolve}
The stationary completion tensor satisfies exactly the same substrate solve as in the redesign, anchoring, action, and kernel-origin notes:
\begin{equation}
\SubReducedEin(\CompLift)_{AB}
=
-\ResidualLift_{AB},
\qquad
\FreeProj\CompLift=0.
\label{eq:subsolve}
\end{equation}
\end{corollary}

\begin{proof}
Project \eqref{eq:Heq} onto \(\CompFree\). Proposition \ref{prop:stiffzero} gives \(\FreeProj\CompLift=0\). Substituting back into \eqref{eq:Heq} yields \(\SubReducedEin(\CompLift)=-\ResidualLift\).
\end{proof}

\begin{remark}
The deeper-origin functional \eqref{eq:micfunc} and the selected constrained auxiliary sector of the kernel-origin note are therefore on-shell equivalent at current scope. The multiplier form of the earlier note is the rigid constrained representative of the same equilibrium rule that appears here as positive free-mode stiffness.
\end{remark}

\section{Observer Pullback and Preservation of the One-Metric GR Package}

\begin{definition}[Completed bridge and observer metrics]
\label{def:completed}
Define
\begin{equation}
\BridgeMetric_{AB}^{\sharp}
\equiv
\BridgeMetric_{AB}
+\BaseLift_{AB}
+\ResponseLift_{AB}
+\CompLift_{AB},
\label{eq:bridgecompleted}
\end{equation}
\begin{equation}
\CompShift_{\mu\nu}
\equiv
\ObserverMap^*\CompLift_{AB},
\qquad
\CompletedMetric
\equiv
\ObserverMetric+\BaseShift+\ResponseShift+\CompShift.
\label{eq:defsigma}
\end{equation}
\end{definition}

\begin{theorem}[The same observer completion solve is recovered]
\label{thm:observer}
The deeper-origin stationary completion tensor satisfies exactly the same observer solve:
\begin{equation}
\ReducedEin(\CompShift)_{\mu\nu}
=
-\ResidualCore{}_{\mu\nu}.
\label{eq:observereq}
\end{equation}
Equivalently,
\begin{equation}
\EinLin(\CompShift)_{\mu\nu}
=
-\ResidualCore{}_{\mu\nu}
+\GaugeBook(\CompShift)_{\mu\nu}.
\label{eq:linsolve}
\end{equation}
\end{theorem}

\begin{proof}
Apply \(\ObserverMap^*\) to \eqref{eq:subsolve}. Using \eqref{eq:intertwine2} and \eqref{eq:liftedsource},
\[
\ReducedEin(\CompShift)
=
\ObserverMap^*\bigl(\SubReducedEin(\CompLift)\bigr)
=
-\ObserverMap^*(\ResidualLift)
=
-\ResidualCore.
\]
This is \eqref{eq:observereq}, and \eqref{eq:linsolve} is the equivalent identity rewritten using \(\ReducedEin=\EinLin-\GaugeBook\).
\end{proof}

\begin{theorem}[The former genuine quartic residual sector remains removed]
\label{thm:removed}
The exact observer equation on the deeper-origin completed branch is
\begin{equation}
G_{\mu\nu}[\CompletedMetric]
=
8\pi G_N\,T_{\mu\nu}^{\mathrm{matter}}[\CompletedMetric,\psi]
+\GaugeBook(\CompShift)_{\mu\nu}
+\CompDec,
\label{eq:newobs}
\end{equation}
where
\begin{equation}
\CompDec
\equiv
\ResidualDec+\CompMix+\CompQuad-8\pi G_N\,\DeltaTComp,
\label{eq:compdec}
\end{equation}
with
\begin{equation}
\DeltaTComp
\equiv
T_{\mu\nu}^{\mathrm{matter}}[\CompletedMetric,\psi]
-T_{\mu\nu}^{\mathrm{matter}}[\CandidateMetric,\psi],
\label{eq:deltat}
\end{equation}
\begin{align}
\CompMix
&\equiv
\EinQuad[\BaseShift+\ResponseShift+\CompShift]
-\EinQuad[\BaseShift+\ResponseShift]
-\EinQuad[\CompShift],
\label{eq:compmix}
\\
\CompQuad
&\equiv
\EinQuad[\CompShift].
\label{eq:compquad}
\end{align}
In particular, the old genuine quartic residual tensor \(\ResidualCore\) is absent from the exact observer equation.
\end{theorem}

\begin{proof}
Expand the Einstein tensor about \(\CandidateMetric\):
\[
G_{\mu\nu}[\CompletedMetric]
=
G_{\mu\nu}[\CandidateMetric]
+\EinLin(\CompShift)_{\mu\nu}
+\CompMix+\CompQuad.
\]
Substitute the fixed candidate equation together with \eqref{eq:linsolve}. The \(\ResidualCore\) terms cancel exactly, leaving
\[
G_{\mu\nu}[\CompletedMetric]
=
8\pi G_N\,T_{\mu\nu}^{\mathrm{matter}}[\CandidateMetric,\psi]
+\GaugeBook(\CompShift)_{\mu\nu}
+\ResidualDec+\CompMix+\CompQuad.
\]
Replacing the candidate matter tensor by
\[
T_{\mu\nu}^{\mathrm{matter}}[\CandidateMetric,\psi]
=
T_{\mu\nu}^{\mathrm{matter}}[\CompletedMetric,\psi]-\DeltaTComp
\]
gives \eqref{eq:newobs}.
\end{proof}

\begin{proposition}[The same one operative metric is preserved]
\label{prop:onemetric}
Matter and light still couple only through the single operative metric \(\CompletedMetric\).
\end{proposition}

\begin{proof}
The present note does not alter the fixed matter route \(\MatterAct[\ObserverMap^*\BridgeMetric,\psi]\). On the completed branch the bridge tensor entering that route is \(\BridgeMetric^\sharp\), whose observer pullback is \(\CompletedMetric\) by Definition \ref{def:completed}. No second matter metric and no direct unsuppressed matter coupling to the auxiliary completion reservoir are introduced.
\end{proof}

\begin{proposition}[The recorded Phase D / E package is preserved]
\label{prop:phaseDE}
The deeper-origin completion preserves the recorded Phase D infrared-spectrum verdict and the Phase E universal-coupling verdict.
\end{proposition}

\begin{proof}
Linearizing \eqref{eq:subsolve} on the admissible homogeneous benchmark background gives
\[
\SubReducedEin(\delta\CompLift)=0,
\qquad
\FreeProj(\delta\CompLift)=0.
\]
The only admissible homogeneous completion datum is the free mode extracted by \(\FreeProj\), so \(\delta\CompLift=0\), hence \(\delta\CompShift=0\). Therefore the linearized observer equation is unchanged relative to the already-screened charge-polarization candidate. No new unsuppressed scalar, vector, ghost, tachyonic, preferred-frame, or second-cone pole appears, and Proposition \ref{prop:onemetric} preserves the same single-metric universal matter route.
\end{proof}

\begin{proposition}[Infrared size is unchanged]
\label{prop:size}
On every controlled infrared family,
\begin{equation}
\CompShift=\mathcal O(\epsIR^4),
\qquad
\GaugeBook(\CompShift)=\mathcal O(\epsIR^4),
\label{eq:ircomp}
\end{equation}
and
\begin{equation}
\CompMix=\mathcal O(\epsIR^6),
\qquad
\CompQuad=\mathcal O(\epsIR^8),
\qquad
\DeltaTComp=\mathcal O(\epsIR^4),
\qquad
\CompDec=\mathcal O(\epsIR^4).
\label{eq:irsizes}
\end{equation}
\end{proposition}

\begin{proof}
By \eqref{eq:irorders}, the source satisfies \(\ResidualLift=\mathcal O(\epsIR^4)\). The deeper-origin solve \eqref{eq:subsolve} is the same reduced solve already used in the redesign, anchoring, action, and kernel-origin notes, imposed in the same admissible class and with the same zero-free-mode condition. Hence \(\CompLift=\mathcal O(\epsIR^4)\), and therefore \(\CompShift=\mathcal O(\epsIR^4)\). The remaining orders follow exactly as in the earlier completion notes from \(\BaseShift=\mathcal O(\epsIR^2)\), \(\ResponseShift=\mathcal O(\epsIR^4)\), and \(\CompShift=\mathcal O(\epsIR^4)\).
\end{proof}

\begin{theorem}[The recorded Phase F controlled GR limit is preserved]
\label{thm:phaseF}
On every controlled infrared family,
\begin{equation}
G_{\mu\nu}[\CompletedMetric_{\epsIR}]
=
8\pi G_N\,T_{\mu\nu}^{\mathrm{matter}}[\CompletedMetric_{\epsIR},\psi_{\epsIR}]
+\GaugeBook(\CompShift_{\epsIR})_{\mu\nu}
+\epsIR^4\mathfrak D_{\mu\nu}^{(\epsIR)},
\label{eq:phaseFnew}
\end{equation}
with \(\mathfrak D_{\mu\nu}^{(\epsIR)}\) uniformly bounded. Thus the same one-metric observer equation reduces to Einstein plus ordinary matter as \(\epsIR\to0\).
\end{theorem}

\begin{proof}
The exact observer equation \eqref{eq:newobs} holds on the deeper-origin completed branch. Proposition \ref{prop:size} gives \(\CompDec=\mathcal O(\epsIR^4)\), so \(\CompDec=\epsIR^4\mathfrak D^{(\epsIR)}\) with \(\mathfrak D^{(\epsIR)}\) uniformly bounded. The explicit remainder \(\GaugeBook(\CompShift)\) remains gauge bookkeeping rather than a genuine observer-side deviation sector, and Proposition \ref{prop:onemetric} preserves the same one operative metric. Therefore the same controlled GR-limit conclusion survives unchanged.
\end{proof}

\section{Claim-Boundary Verdict}

\begin{theorem}[Primary deeper-origin verdict]
\label{thm:verdict}
At the disciplined scope of the present note, the positive result is exactly the following.
\begin{enumerate}
    \item The kernel-origin note's three current-scope selection principles now themselves receive a still-deeper microscopic-equilibrium origin.
    \item Branch-preserving quadratic response arises because the unsourced completion branch is stationary, so its first nontrivial response is the Hessian of the same local bridge geometry; within the fixed branch data and derivative order, that Hessian is the already-recorded bridge response \(\SubEinLin\).
    \item Reduced-gauge neutrality is forced because microscopic completion weights are neutral under reduced-gauge relabeling and therefore depend only on the neutral combination \(\GaugeAux-\SubGaugeVec(H)\), which yields the same reduced operator \(\SubReducedEin\).
    \item Equilibrium suppression of source-free homogeneous data is forced because the free homogeneous sector carries positive stiffness while the fixed lifted quartic source carries no free component; therefore \(\FreeProj\CompLift=0\).
    \item The same substrate and observer solves are recovered:
    \[
    \SubReducedEin(\CompLift)=-\ResidualLift,
    \qquad
    \ReducedEin(\CompShift)=-\ResidualCore.
    \]
    \item The same one operative metric, the same Phase D / E package, the same Phase F controlled GR limit, and the same removal of the former genuine quartic residual tensor are preserved.
    \item Stronger promotion language is still not justified here, because the note supplies a deeper microscopic origin of the present selection principle without claiming a final ultraviolet completion or any theorem broader than the current one-metric claim boundary.
\end{enumerate}
\end{theorem}

\begin{proof}
Item (1) is Theorem \ref{thm:principles}. Item (2) is Proposition \ref{prop:quadratic} and Proposition \ref{prop:hessian}. Item (3) is Proposition \ref{prop:neutralcombo}, Proposition \ref{prop:reduced}, and Corollary \ref{cor:neutral}. Item (4) is Proposition \ref{prop:stiffzero} and Corollary \ref{cor:projector}. Item (5) is Corollary \ref{cor:subsolve} and Theorem \ref{thm:observer}. Item (6) is Theorem \ref{thm:removed}, Proposition \ref{prop:onemetric}, Proposition \ref{prop:phaseDE}, and Theorem \ref{thm:phaseF}. Item (7) is the claim-boundary discipline stated in the introduction.
\end{proof}

\begin{corollary}[What remains next]
The primary active burden is now narrower again. It is no longer to justify the current-scope selection principle itself, since that deeper origin is now recorded at microscopic-equilibrium scope. The next honest burden is either a still deeper primitive origin of that equilibrium reservoir, symmetry, or stiffness data, or another comparably concrete observer-equation gain on the same one-metric GR line. The anchored positive-pair continuation remains side work unless it directly supplies one of those.
\end{corollary}

\section{Conclusion}

The kernel-origin note already showed that the completion kernel, reduced-gauge auxiliary block, and free-mode projector are selected at disciplined current-scope level by branch-preserving quadratic response, reduced-gauge neutrality, and equilibrium suppression of source-free homogeneous data. The present manuscript moves one layer deeper again. It does not write another same-level completion action. Instead it explains why those three selection principles arise from a deeper microscopic-equilibrium completion reservoir on the fixed one-metric charge-polarization branch.

At the unsourced reference branch, stationarity makes the first nontrivial completion response quadratic. Local branch covariance then forces that quadratic response to be the same bridge Einstein--Hilbert Hessian already built into the active branch rather than a new principal operator. Reduced-gauge relabeling neutrality forces the completion reservoir to depend only on the neutral combination \(\GaugeAux-\SubGaugeVec(H)\), which reproduces the same reduced operator \(\SubReducedEin=\SubEinLin-\SubGaugeBook\). Positive equilibrium stiffness of the source-free homogeneous sector, together with the fact that the fixed lifted quartic source has no overlap with that sector, forces the zero-free-mode condition \(\FreeProj\CompLift=0\).

Consequently the same completion solve is recovered again:
\[
\SubReducedEin(\CompLift)=-\ResidualLift,
\qquad
\ReducedEin(\CompShift)=-\ResidualCore.
\]
The former genuine quartic residual tensor therefore remains absent from the exact observer equation, the one operative metric remains preserved, and the recorded Phase D / E and Phase F package stays intact.

The scope remains disciplined. This note does not claim a final ultraviolet completion, does not widen the current finite-amplitude theorem boundary, and does not license stronger promotion language. Its positive result is narrower and exactly the one now needed: the selection principle used to choose the completion kernel, projector, and constrained auxiliary sector now itself has a still-deeper microscopic origin on the same one-metric GR line, while the next open burden is pushed below that layer again.

\end{document}
